\chapter{Glossaire}
\label{ch:glossary}

\begin{description}

  \item[Administrateur]\index{administrateur} Compte utilisateur
        avec le rôle \texttt{admin}. Le seul rôle qui peut ouvrir
        le panneau \texttt{/admin} et éditer les utilisateurs, les
        données de référence, la corbeille ou les paramètres
        système.

  \item[Anonyme]\index{anonyme} Visiteur sans cookie de session.
        Accès en lecture seule au catalogue public. Pas de
        catalogage, pas de prêt, pas d'édition.

  \item[Audit]\index{audit} Le rituel « je parcours les étagères
        et je confirme que tout est à la place que la base
        indique » décrit en section~\ref{sec:audit}.

  \item[BDGest]\index{BDGest} Catalogue BD francophone utilisé comme
        source de métadonnées pour ce type de média.

  \item[BnF]\index{BnF} \textit{Bibliothèque nationale de France}.
        Catalogue de la BNF, utilisé comme source principale de
        métadonnées pour les livres en français.

  \item[Catalogage] L'acte d'enregistrer un nouveau titre et/ou
        volume dans la base. Typiquement piloté au code-barres :
        scanner, confirmer, coller l'étiquette, recommencer.

  \item[Code L]\index{code L} Étiquette code-barres d'un lieu de
        stockage. Format : \texttt{L} suivi d'un numéro
        zéro-paddé.

  \item[Code V]\index{code V} Étiquette code-barres d'un volume.
        Format : \texttt{V} suivi d'un numéro zéro-paddé.

  \item[Couverture]\index{couvertures} La face visuelle avant d'un
        livre / disque / vinyle, téléchargée depuis un fournisseur
        de métadonnées et mise en cache sous \texttt{COVERS\_DIR}
        sur disque.

  \item[CSP] Content Security Policy. En-tête imposé par le
        navigateur qui restreint quels scripts et styles la page
        peut exécuter. mybibli applique une politique stricte qui
        interdit scripts et styles inline.

  \item[CSRF (jeton)]\index{CSRF} Secret par session intégré dans
        chaque formulaire. mybibli rejette les requêtes POST /
        PUT / DELETE dont le jeton ne correspond pas — défense
        contre les forgeries de requête inter-sites.

  \item[Dewey (code)]\index{Dewey} Le numéro de la Classification
        décimale de Dewey, utilisé pour la navigation par sujet
        du catalogue. Optionnel par titre.

  \item[EAN-13] Le standard de code-barres à 13 chiffres. Les ISBN
        sont encodés en EAN-13 depuis 2007 ; les CD, DVD et autres
        supports portent aussi des codes EAN-13.

  \item[Emprunteur]\index{emprunteur} Personne qui emprunte des
        volumes à la bibliothèque. Modélisé comme sa propre entité
        (nom, contact optionnel, historique de prêts).

  \item[Exemplaire] Synonyme de \emph{volume} (l'objet physique).
        Utilisé par certaines bibliothèques françaises ; mybibli
        préfère \emph{volume}.

  \item[État de volume]\index{état de volume} Ligne de données de
        référence décrivant l'état d'un volume : Neuf, Bon, Usé,
        Endommagé, \dots.

  \item[Garde du dernier admin actif]\index{garde du dernier admin}
        Règle qui empêche de désactiver ou supprimer définitivement
        le dernier utilisateur de rôle \texttt{admin}. Voir
        section~\ref{sec:users}.

  \item[Genre]\index{genres} Classification libre du sujet d'un
        titre (Fiction, Documentaire, Polar, \dots). Définie dans
        la table de référence des genres.

  \item[Hard-delete] Suppression d'une ligne pour de bon. mybibli
        ne hard-delete que depuis le panneau \texttt{/admin >
        Trash} et la tâche d'auto-purge en arrière-plan ; tout le
        reste est un soft-delete.

  \item[HTMX] Bibliothèque JavaScript qu'utilise mybibli pour
        échanger des fragments de page sans rechargement complet.
        Les exploitants n'ont pas besoin de comprendre HTMX pour
        utiliser mybibli.

  \item[ISBN] International Standard Book Number. Identifiant
        imprimé au dos des livres, encodé en code-barres EAN-13.

  \item[Lieu]\index{lieu} L'endroit physique où vit un volume.
        Modélisé comme un arbre (pièce → étagère → rangée),
        identifié par un code L.

  \item[MariaDB] Le moteur de base où mybibli stocke ses données.
        Fork de MySQL.

  \item[Multi-position (volume)]\index{multi-position} Un volume
        physique unique qui contient plusieurs positions dans une
        série — typiquement un album BD omnibus.

  \item[MusicBrainz]\index{MusicBrainz} Catalogue de métadonnées
        musicales maintenu par des bénévoles, utilisé comme source
        principale pour les médias audio.

  \item[OMDb]\index{OMDb} Open Movie Database. Source de
        métadonnées pour films et séries TV.

  \item[Open Library]\index{Open Library} Catalogue ouvert de
        livres porté par des bénévoles, utilisé comme source de
        métadonnées.

  \item[Optimistic locking] Technique qu'utilise mybibli pour
        détecter les éditions concurrentes. Deux utilisateurs qui
        éditent le même titre obtiennent une erreur claire à la
        seconde sauvegarde au lieu de s'écraser silencieusement.

  \item[Prêt]\index{prêt} Enregistrement qu'un volume précis est
        actuellement sorti chez un emprunteur précis. Note la date
        de prêt, la date de retour prévue, la date de retour, et
        le lieu pré-prêt du volume.

  \item[Réactivation] Restauration d'une ligne soft-deleted. Pour
        les utilisateurs, efface \texttt{deleted\_at} et leur
        permet de se reconnecter. Pour les données de référence,
        le terme couvre aussi le fait de re-créer une ligne avec
        le même nom qu'une soft-deleted, ce qui défait la
        suppression de manière transparente.

  \item[Série]\index{série} Suite ordonnée et nommée de titres
        liés (par exemple une trilogie ou une série BD en cours).
        La position de chaque titre dans la série est un petit
        entier ou une plage (pour les omnibus).

  \item[Setup wizard]\index{setup wizard} Le flux en quatre étapes
        à \texttt{/setup} qu'une instance vide présente jusqu'à
        ce qu'un compte admin existe.

  \item[Soft-delete]\index{soft-delete} Positionnement de la
        colonne \texttt{deleted\_at} d'une ligne à « maintenant »
        au lieu de supprimer physiquement la ligne. Les lignes
        soft-deleted sont invisibles aux requêtes normales mais
        apparaissent dans \texttt{/admin > Trash} pendant
        30\,jours, après quoi l'auto-purge les hard-delete.

  \item[Titre] L'entité bibliographique (un livre, un film, un
        album). Un titre peut avoir plusieurs volumes — par
        exemple, deux exemplaires du même roman sur des étagères
        différentes.

  \item[TMDb]\index{TMDb} The Movie Database. Source de
        métadonnées communautaire pour films et séries TV.

  \item[Type de média]\index{types de médias} La nature de l'objet
        que représente un titre : Livre, BD / comic, Audio, Film /
        série. Pilote la sélection des fournisseurs et l'affichage
        des champs.

  \item[Type de nœud]\index{type de nœud} Libellé appliqué aux
        lieux de stockage pour indiquer quel type de contenant
        ils sont (pièce, étagère, rangée, boîte). Pur libellé —
        aucune contrainte sémantique.

  \item[Volume]\index{volume} L'objet physique sur l'étagère — un
        exemplaire d'un titre. Un titre peut avoir plusieurs
        volumes ; chacun porte son propre code V, son lieu et son
        état.

\end{description}
